\documentclass{article}


\usepackage{arxiv}
\usepackage{enumitem}
\usepackage[utf8]{inputenc} % allow utf-8 input
\usepackage[T1]{fontenc}    % use 8-bit T1 fonts
\usepackage{hyperref}       % hyperlinks
\usepackage{url}            % simple URL typesetting
\usepackage{booktabs}       % professional-quality tables
\usepackage{amsfonts}       % blackboard math symbols
\usepackage{nicefrac}       % compact symbols for 1/2, etc.
\usepackage{microtype}      % microtypography
\usepackage{lipsum}
\usepackage{fancyhdr}       % header
\usepackage{graphicx}       % graphics
\usepackage{amsmath}
\usepackage{lyu}
\usepackage{amsthm}
\usepackage{amssymb}
\usepackage{xcolor}
\newtheorem{theorem}{Theorem}[section]
\newtheorem{corollary}{Corollary}[theorem]
\newtheorem{lemma}[theorem]{Lemma}
\theoremstyle{definition}
\newtheorem{definition}{Definition}[section]
\newtheorem{assumption}{Assumption}
\newtheorem{proposition}[theorem]{Proposition}
\newtheorem{remark}[theorem]{Remark}

\pagestyle{fancy}
\thispagestyle{empty}
% \rhead{ \textit{ }} 

% Update your Headers here
% \fancyhead[LO]{}
% 
% \fancyhead[RE]{Firstauthor and Secondauthor} % Firstauthor et al. if more than 2 - must use \documentclass[twoside]{article}


% Update your Headers here
\fancyhead[LO]{Multiscale Nudge}




  
\title{
Multiscale Nudge: Cross-scale Data Assimilation from Macroscopic Observations to Microscopic Dynamics
}

\author{
  Liyao Lyu\\
  Department of Mathematics \\
  University of California, Los Angeles \\
  Los Angeles, CA 90024, USA\\
  \texttt{lyuliyao@math.ucla.edu} \\
  \AND
   Xinyue Yu \\
  Department of Mathematics \\
  University of California, Los Angeles \\
  Los Angeles, CA 90024, USA\\
  \texttt{tracy@math.ucla.edu} \\
  \AND
   Hayden Schaeffer  \\
  Department of Mathematics \\
  University of California, Los Angeles \\
  Los Angeles, CA 90024, USA\\
  \texttt{hayden@math.ucla.edu} \\
  %% Affiliation \\
  %% Address \\
  %% \texttt{email} \\
  %% \And
  %% Coauthor \\
  %% Affiliation \\
  %% Address \\
  %% \texttt{email} \\
  %% \And
  %% Coauthor \\
  %% Affiliation \\
  %% Address \\
  %% \texttt{email} \\
}


\begin{document}

Consider the Vlasov--Poisson system in characteristic form,
\begin{align}
    \mathrm{d} X_t^i &= V_t^i\,\mathrm{d} t,\\
    \mathrm{d} V_t^i &= E(X_t^i,t)\,\mathrm{d} t,
\end{align}
with self-consistent electric field $E(x,t) = -\nabla_x\phi(x,t)$, where the potential solves the Poisson equation
\begin{equation}\label{equ:poisson}
    -\Delta_x\phi(x,t) = \rho(x,t) - 1, 
    \qquad 
    \rho(x,t) = \int f(x,v,t)\,\mathrm{d} v.
\end{equation}


Suppose we only observe the potential, $y(x,t) := \phi(x,t)$, and wish to reconstruct the full phase-space density $f(x,v,t)$. This inverse problem is severely ill-posed: $\phi$ depends on $f$ only through its spatial marginal $\rho$, so infinitely many distributions produce the same observation. To extract velocity information, we use multiple historical snapshots and nudge the state through the cost functional
\begin{equation}
    J_k(g) = \frac{1}{2}\int_0^{T_w} K(s)\,\bigl\|H(B_s g) - y_k(s)\bigr\|^2\,\mathrm{d} s,
\end{equation}
where
\begin{itemize}
    \item $H : f \mapsto \phi[f]$ is the (linear) observation operator, defined by solving \eqref{equ:poisson};
    \item $B_s$ is the \emph{backward evolution operator}, $B_s g := f_g(s)$, with $f_g$ solving
    \begin{equation}\label{equ:bk_operator}
        \partial_s f_g - v\cdot\nabla_x f_g - E[f_g]\cdot\nabla_v f_g = 0, 
        \qquad 
        f_g(0) = g;
    \end{equation}
    \item $K(s)$ is a temporal weight and $y_k(s)$ is the $k$-th observed snapshot.
\end{itemize}


A formal variational calculation gives
\begin{equation}
    \delta J_k(g)[\eta] 
    = \int_0^{T_w} K(s)\, B_s'(g)^{*}\, H^{*}\bigl(H(B_s g) - y_k(s)\bigr)\,\mathrm{d} s,
\end{equation}
where $B_s'(g)^{*}$ and $H^{*}$ are the adjoints of the linearized backward operator and the observation operator, respectively. Since $H$ is linear, $H'(f) = H$; the starred operators are the adjoints used to push the residual back onto the initial-condition space.


To identify $B_s'(g)$, perturb the initial datum and set
\begin{equation}
    h(s) := \left.\frac{\mathrm{d}}{\mathrm{d}\epsilon}\, f_{g+\epsilon\eta}(s)\right|_{\epsilon=0} 
    = B_s'(g)\,\eta.
\end{equation}
Differentiating \eqref{equ:bk_operator} in $\epsilon$ yields the linearized Vlasov equation
\begin{equation}
    \partial_s h - v\cdot\nabla_x h - E[f_g]\cdot\nabla_v h - \delta E[h]\cdot\nabla_v f_g = 0, 
    \qquad 
    h(0) = \eta,
\end{equation}
where the field perturbation $\delta E[h] = -\nabla_x\delta\phi[h]$ is determined by
\begin{equation}
    -\Delta_x\delta\phi[h] = \rho_h, 
    \qquad 
    \rho_h(x,s) = \int h(x,v,s)\,\mathrm{d} v.
\end{equation}
The additional term $\delta E[h]\cdot\nabla_v f_g$ reflects the self-consistent nature of the field: a perturbation of $f$ perturbs $\rho$, hence $\phi$, hence the transport field itself.

\bibliographystyle{unsrt}  

\appendix
\end{document}
